\scanpage{522}
\mspage{433}
% B01
$c\Longrightarrow a$, \deleted{\unsure{OPS}} on peut donc supposer que $C'=C$, $C'\longrightarrow C$ le foncteur identique, et prouver que \unsure{si} c) est \unsure{satisfait|satisfaite} et si $C$ est petite et sans projecteurs, alors $E$ est accessible et $f$ accessible. Or c'est ce qu'affirme 4.15.1.

% B02
\uline{\unsure{Remarque|Remarques}} 4.15.5. Si $f:E\longrightarrow C$ est un foncteur pl fidèle, il est immédiat que la condition $H$ est satisfaite (car $E_\lambda\to\unsure{\mathbb{\Delta}|\Delta}^1$ sera \unsure{plein\textsuperscript{t}|plein\textsuperscript{t}.} fidèle, donc $E_\lambda$ est équivalente à $e$ ou à $\mathbb{\Delta}^1$, donc accessible, et tout foncteur de $E_\lambda$ dans une cat. quelconque est accessible \unsure{. --|...|..}).
% B03
Mais si on suppose $C$ petite, \uline{sans} supposer $C$ sans projecteurs, il ne s'ensuit pas que $E$ soit accessible, \deleted{\unsure{ou}} ou ce qui revient au $=$ (étant \deleted{\unsure{p}}ess. petite), que $E$ soit \unsure{Karoubienne|karoubienne}. \deleted{\unsure{Prenons}} (L'exemple le plus trivial est celui où $E=C$, $f=\mathrm{id}_C$, $E$ n'est karoubienne que si \deleted{\unsure{E}} $C$ l'est !)

% B04
Mais $=$ en supposant $C$ \deleted{\unsure{kar}} une petite catégorie karoubienne, il ne s'ensuit pas forcément que $E$ l'est. \deleted{\unsure{en}} En effet, \uline{tte} petite catégorie $E$ (karoubienne ou non) se plonge comme sous-catégorie pleine dans \deleted{une} son enveloppe karoubienne $C$ -- il suffit donc de prendre $E$ non karoubienne, L'exemple

% archive sheet number: 521
